\section{集合的标度。结构}

\begin{enumerate}
\def\labelenumi{\arabic{enumi}.}
\tightlist
\item
  例如，给定三个不同的集合 E、F、G，我们可以通过取它们的子集集合，或通过取其中一个与自身的乘积，或再次通过按一定顺序取其中两个的乘积，从它们构造出其他集合。这样我们得到十二个新集合。如果我们将这些新集合添加到三个原始集合 E、F、G 中，我们可以对这十五个集合重复同样的操作，省略那些给出已得到集合的操作；依此类推。一般地，通过此过程（按照明确的方案）得到的任何集合，都称为以 E、F、G 为基的集合标度中的集合。
\end{enumerate}

例如，设M、N、P为该尺度上的三个集合，并设 \(R\{x, y, z\}\) 是 M、N、P 的典型元素 x、y、z 之间的关系。则 R 定义了 \(M \times N \times P\) 的一个子集，因此（通过典范对应）是 \((\mathbf{M} \times \mathbf{N}) \times \mathbf{P}\) 的子集，即 \(\mathfrak{P}((\mathbf{M} \times \mathbf{N}) \times \mathbf{P})\) 的一个元素。因此，给出同一尺度中若干集合的元素之间的关系，等同于给出该尺度中另一个集合的一个元素。类似地，给出从 M 到 N 的映射，例如，等同于（通过考虑该映射的图像）给出 \(M \times N\) 的子集，即 \(\mathfrak{P}(\mathbf{M} \times \mathbf{N})\) 的一个元素，这又是该尺度中的一个集合。最后，给出（例如）M 的两个元素，相当于给出乘积集合 \(M \times M\) 中的一个元素。

因此，给定尺度中若干集合的某些元素、这些集合的泛元素之间的关系，以及其中某些集合的子集到其他集合的映射，归根结底都归结为给定该尺度中某个集合的一个元素。

\begin{enumerate}
\def\labelenumi{\arabic{enumi}.}
\setcounter{enumi}{1}
\tightlist
\item
  我们前面（§6）说过，集合 \(\mathfrak{P}(\mathbf{E} \times \mathbf{E})\) 的一个元素 C 在具有性质 (a) \(C \circ C \subset C\) 和 (b) 时，在 E 上定义了一个序结构。 \(C \cap \overline{C} = \Delta\) .
\end{enumerate}

一般而言，考虑一个集合尺度中的集合M，其基为举例而言由三个集合E、F、G组成。我们给自己若干明确陈述的关于M的典型元素的性质，并设T为这些性质所定义的M的子集的交集。一个元素 \(\sigma\) T 被称为在 E、F、G 上定义 T 物种的一个结构。因此，T 物种的结构由从 E、F、G 形成 M 的构造方案以及定义 T 的性质（这些性质被称为这些结构的公理）来刻画。我们对同一物种的所有结构赋予一个特定名称。每一个由命题“ \(\sigma \in T\) ''（即定义T的公理）被称为属于T类结构理论；例如，§6中陈述的命题属于有序集结构理论。

在最后一个例子中，公理可以针对完全任意的基集E来陈述。因此，我们给满足这些公理的结构赋予相同的名称，而不论它们定义在哪个集合上；从这些公理推导出的命题在任何集合中都成立，因为它们的表述不涉及集合E的任何特殊性质。每当公理具有这种性质时，这些评论都适用（*）。

最常情况下，当标度与由若干集合E、F、G构成的基一起使用时，这些集合中的一个，比如E，在所考虑的结构中起着主导作用。因此，出于语言上的滥用，这些结构被称为定义在集合E中，而F和G被视为辅助集合。

最后，为了简化语言，通常会给一个被赋予了特定种类结构的集合一个专门的名称。因此我们谈论有序集，在本系列后续部分中，我们将定义群、环、域、拓扑空间、一致空间等概念，所有这些词都表示赋予了某种结构的集合。

\begin{enumerate}
\def\labelenumi{\arabic{enumi}.}
\setcounter{enumi}{2}
\item
  考虑同一物种T的结构，其中T是集合M在集合尺度中的一个子集。如果我们对定义T的公理附加新的“公理”，这样得到的公理系统就定义了M的一个子集U，且U包含于T。物种U的结构被称为比物种T的结构更丰富。例如，全序集的结构比有序集的结构更丰富，因为元素C属于 \(\mathfrak{P}(\mathbf{E} \times \mathbf{E})\) 定义这种结构的公理必须满足额外的公理。 \(C \cup C^{-1} = E \times E\) .
\item
  设M、M'为同一尺度下的两个集合，设以E、F、G为基。设T为M的子集，T'为M'的子集，各自由某些明确陈述的公理定义。每当我们能明确地定义T到T'的一一映射时，我们考虑两个元素 \(\sigma \in T\) , \(\sigma' \in T'\) 在此映射下彼此对应的这些对象，被视为在E、F、G上定义了相同的结构；而定义T与T'的公理系统则被称为等价的。
\end{enumerate}

拓扑结构提供了这种情况的一个例子；它们可以通过若干等价的公理系统来定义，其中有两个系统特别有用（见《一般拓扑学》第一章第1节）。

\begin{enumerate}
\def\labelenumi{\arabic{enumi}.}
\setcounter{enumi}{4}
\tightlist
\item
  设E、F、G为三个集合，并假设我们给出了从E、F、G到另外三个集合的双射映射。 \(E'\) , \(F'\) , \(G'\) 分别地。由于我们知道如何定义双射映射到子集集合（§2，第9节）和乘积集合（§3，第14节）的典范扩张，我们可以逐步定义给定双射到两个集合M、M'的扩张，这两个集合分别按照基于E、F、G和基于E'、F'、G'的集合尺度中的同一方案构造。设f为由此得到的从M到M'的双射。如果σ是E、F、G上的一个结构，它是M的子集T的一个元素，我们说 \(f(\sigma)\) 是通过给定的E到E'、F到F'、G到G'的双射将结构σ搬运到E'、F'、G'上得到的结构。关于E、F、G上的结构σ的每一个命题（通过使用适当的扩张）都会产生一个关于 \(f(\sigma)\) 上的结构E'、F'、G'的命题。
\end{enumerate}

反之，一个结构 \(\sigma\) 在 E、F、G 上以及一个结构 \(\sigma'\) 在 E'、F'、G' 上，若 \(\sigma'\) 可以通过搬运 \(\sigma\) 分别通过 E、F、G 到 E'、F'、G' 的双射来获得，则称它们是同构的；这些映射随后被称为构成 \(\sigma\) 到 \(\sigma'\) 上的同构.

当我们关注单个集合 E 上的结构时，E 到 \(E'\) 的双射（它将 \(\sigma\) 搬运到 \(\sigma'\) ）也称为集合E（赋予结构 \(\sigma\) ）到集合 \(E'\) （赋予结构 \(\sigma'\) ）上的同构.

当F和G是两个辅助集合，且这两个集合的双射是F和G到自身的恒等映射时，这个映射也称为同构。

赋予了结构 \(\sigma\) 的集合E到自身的同构称为自同构。

当存在一个集合E的同构，该集合E被赋予一种结构时 \(\sigma\) 到集合 \(E'\) 的结构，配备有结构 \(\sigma'\) 上时，通常方便将 E 等同于 \(E'\) ，即，在基于E的集合标架中，给集合M的一个元素与它在f到集合M的适当扩张下的像赋予相同的名称。

\begin{enumerate}
\def\labelenumi{\arabic{enumi}.}
\setcounter{enumi}{5}
\item
  给定一组公理，在集合标架中定义一个集合M的子集T，在谈论满足这些公理的结构之前，我们应当确保集合T不一定是空的；否则这些公理就被称为矛盾的或不一致的。
\item
  可能发生这样的情况：一组在集合上定义结构的公理可以对任意集合陈述，但是当我们考虑满足这些公理且定义在两个不同集合E、F上的两个结构时，从公理中我们发现这些结构（如果存在的话）必然是同构的（这特别意味着E和F是等势的）。那么满足这些公理的结构的理论被称为单价的；否则它们被称为多价的。
\end{enumerate}

整数理论、实数理论和经典欧几里得几何是单价理论；有序集合理论、群论和拓扑学是多价理论。多价理论的研究是现代数学区别于经典数学的最显著特征。

